\label{sec:datasets}

The latest work plan separates the local railway data into two complementary
datasets. Domain-QA emphasizes bilingual terminology and translation.
Domain-RegQA emphasizes evidence-grounded regulatory question answering. This
separation is important because the two datasets test different failure modes:
terminology tasks reward concise domain wording, while regulatory QA rewards
faithfulness to source clauses.

\begin{table*}[t]
\centering
\caption{Dataset statistics used in the planned evaluation and adaptation workflow. Domain-QA is a railway bilingual terminology and translation dataset. Domain-RegQA is an evidence-grounded regulatory QA dataset generated by rule-based extraction rather than by LLM synthesis.}
\label{tab:dataset-statistics}
\resizebox{\textwidth}{!}{%
\begin{tabular}{llrrrrl}
\toprule
Dataset & Split & Samples & Main categories & Evidence field & Gold span & Main purpose \\
\midrule
Domain-QA & train & 22,164 & 4 & no & no & QLoRA adaptation and bilingual/domain QA \\
Domain-QA & validation & 2,770 & 4 & no & no & validation and early comparison \\
Domain-QA & test & 2,770 & 4 & no & no & held-out domain QA evaluation \\
Domain-RegQA & train & 1,600 & 9 & yes & yes & regulatory QA adaptation/ablation \\
Domain-RegQA & validation & 200 & 6 & yes & yes & evidence-grounded validation \\
Domain-RegQA & test & 200 & 7 & yes & yes & evidence-grounded domain evaluation \\
\bottomrule
\end{tabular}%
}
\end{table*}


\subsection{Domain-QA}

Domain-QA is stored under \texttt{data/domain}. It contains 22,164 training
examples, 2,770 validation examples, and 2,770 test examples. The train split
contains 10,153 Chinese-to-English terminology examples, 10,153
English-to-Chinese terminology examples, 929 Chinese-to-English translation
examples, and 929 English-to-Chinese translation examples. The validation and
test splits each contain 1,269 examples for each terminology direction and 116
examples for each translation direction.

Each record contains a prompt, a reference answer, a combined text field,
category information, and source metadata. The answer style is intentionally
short so that exact match, character-level F1, token-level F1, reference
containment, and length ratio can be computed reproducibly. Exact match is
strict and may underestimate semantically correct translations, so
character-level and token-level F1 are treated as complementary domain metrics.

\subsection{Domain-RegQA}

Domain-RegQA is stored under \texttt{data/domain\_regqa}. It contains 1,600
training examples, 200 validation examples, and 200 test examples. The dataset
is evidence-grounded: each sample includes an evidence passage, an answer span,
and generation metadata. The split was generated by rule-based extraction from
source material rather than by LLM synthesis, which reduces the risk that the
evaluation target merely reflects another model's output style.

The training split covers nine regulatory categories, including clause,
requirement, inspection, standard, prohibition, definition, responsibility,
principle, and scope questions. The validation and test splits cover the most
frequent categories with smaller counts. Domain-RegQA is harder than
Domain-QA because an acceptable answer must remain close to a source passage
instead of simply producing a known bilingual term.

\subsection{Dataset Use in the Experiment Plan}

Domain-QA is used for the first completed adapter family, called Adapter-A in
this draft. Domain-RegQA is reserved for Adapter-B. A mixed Domain-QA +
Domain-RegQA configuration is reserved for Adapter-C, while a completion-only
mixed configuration is reserved for Adapter-D. These variants are retained as
placeholders in the current manuscript so that the final paper can compare
whether bilingual terminology data, regulatory evidence data, or mixed
adaptation gives the best trade-off.
